%=====================================================================
%  08-KANAL.TEX — KANAL KUANTUM
%=====================================================================
%  PERBAIKAN dari versi lama:
%    * "p_x Z\rho Z" -> "p_z Z\rho Z" (salah indeks pada kanal Pauli);
%    * "qubt" -> "qubit".
%=====================================================================
\section{CHANNEL}

\begin{frame}
    \frametitle{Kanal Kuantum 1}
    \begin{columns}[T]
        \begin{column}{0.5\textwidth}\vspace{-0.5cm}
            \small\justifying
            \begin{itemize}
                \item The Pauli Channel \\
                    {\fontsize{4}{6}\selectfont\justifying
                    Kanal Pauli mentransformasikan keadaan masukan $\rho$
                    sebagai:
                    \begin{equation}
                        \rho \rightarrow \mathcal{C}_P(\rho)
                        = (1-p)\rho + p_x X\rho X + p_y Y\rho Y + p_z Z\rho Z,
                    \end{equation}
                    dengan $X$, $Y$, $Z$ matriks Pauli satu-qubit
                    \begin{equation}
                        X=\begin{bmatrix}0&1\\1&0\end{bmatrix},\quad
                        Y=\begin{bmatrix}0&-i\\ i&0\end{bmatrix},\quad
                        Z=\begin{bmatrix}1&0\\0&-1\end{bmatrix}.
                    \end{equation}
                    Probabilitas depolarisasi total $p$ adalah jumlah
                    probabilitas galat Pauli:
                    \begin{equation}
                        p=p_X+p_Y+p_Z.
                    \end{equation}
                    Pada waktu sesaat $t$, peluang galat bergantung pada waktu
                    relaksasi $T_1$ dan waktu \textit{dephasing} $T_2$:
                    \begin{equation}
                        \begin{aligned}
                            p_x &= p_y = \frac{1}{4}\left(1-e^{-t/T_2}\right)\\
                            p_z &= \frac{1}{4}\left(1+e^{-t/T_1}-2e^{-t/T_2}\right)
                        \end{aligned}
                    \end{equation}}
            \end{itemize}
        \end{column}

        \begin{column}{0.5\textwidth}\vspace{-0.5cm}
            \small\justifying
            \begin{itemize}
                \item The Depolarizing Channel \\
                    {\fontsize{4}{6}\selectfont\justifying
                    Kanal depolarisasi menyatakan matriks masukan $\rho_i$
                    dapat mengalami depolarisasi: sebagian informasi hilang
                    secara acak. Jika $p=0$ kanal tidak mendepolarisasi;
                    jika $p=1$ kanal menghasilkan keadaan acak penuh
                    (\textit{fully random state}):
                    \begin{equation}
                        \mathcal{N}(\rho_i)=p\,\frac{I}{2}+(1-p)\rho_i.
                    \end{equation}
                    Untuk dua pengguna ortogonal, \textit{mixed input state}
                    \begin{equation}
                        \rho=\sum_i p_i\rho_i=p_0\rho_0+(1-p_0)\rho_1.
                    \end{equation}
                    Setelah transformasi $\mathcal{N}$:
                    \begin{equation}
                        \begin{aligned}
                            \mathcal{N}\!\left(\sum_i p_i\rho_i\right)
                                &= \mathcal{N}\!\left(p_0\rho_0+(1-p_0)\rho_1\right)\\
                                &= p\,\frac{1}{2}I+(1-p)\left(p_0\rho_0+(1-p_0)\rho_1\right).
                        \end{aligned}
                    \end{equation}}
            \end{itemize}
        \end{column}
    \end{columns}
\end{frame}

%---------------------------------------------------------------------
%  Template slide lanjutan (Kanal Kuantum 2) — ganti konten di bawah
%  ini atau duplikat frame di atas untuk menambah slide baru.
%---------------------------------------------------------------------
%\begin{frame}
%    \frametitle{Kanal Kuantum 2}
%    \begin{columns}[T]
%        \begin{column}{0.5\textwidth}
%            \begin{itemize}\small
%                \item Topik berikutnya...
%            \end{itemize}
%        \end{column}
%        \begin{column}{0.5\textwidth}
%            Kolom kanan...
%        \end{column}
%    \end{columns}
%\end{frame}
